\section{Discussion and Policy Implications}\label{sec:discussion}

The deterministic results under the primary channel-based specification leave a narrow policy conclusion: because no cell crosses in the non-stress scenarios (low, mid, frontier), AI-linked fiscal space should be treated as a contingent, long-horizon opportunity rather than as revenue available for the next budget cycle. The framework separates the technological shock from the country--instrument cells in which it could cover a commitment after domestic translation, fiscal capture, and debt consistency. It identifies financing conditions; it does not estimate causal effects or rank the programs' social value.

\subsection{Comparative economic interpretation}\label{subsec:discussion_comparative}

The comparative pattern is governed by the interaction of recurring cost, fiscal conversion, and debt space rather than by a single technology margin. Chilean GMI combines a comparatively low cost with sufficient conversion and debt room only under diagnostic stress; survival of the loaded-cost variant shows that ideal targeting is not doing all the work. Peru's pension illustrates the other mechanism: a low enough cost can permit accounting coverage while the inherited primary-balance gap absorbs the revenue needed for debt consistency. Revenue used to repair that gap cannot simultaneously finance a recurrent entitlement.

The same interaction makes the Colombian and Mexican non-crossings informative. Their closest low-cost cells remain below a common financing boundary, which bounds what the evaluated inputs can support without implying that AI lacks value or that the programs lack merit. Across the non-stress scenarios, the shared mechanism is a joint shortfall: domestic translation at the maintained fiscal architecture is insufficient, and reform-level capture at the maintained technological gain is also insufficient. Translation and capture are complements, while instrument cost and debt space determine how close each cell comes to the boundary.

This interpretation is also what distinguishes the framework from a static redistribution exercise. A static comparison can ask how a given tax base or transfer changes disposable income. The threshold approach asks whether a scenario-contingent fiscal dividend is large and persistent enough to sustain a specific policy under the country's budget constraint. The former remains necessary for welfare design; the latter disciplines when a projected technological dividend may enter a financing argument. The scenario and uncertainty qualifications established in Section~\ref{sec:design} therefore apply once here without changing the comparative policy reading of Section~\ref{sec:results}.

\subsection{Policy implications}\label{subsec:discussion_policy}

The first implication is a financing rule. Under the non-stress scenarios, AI-generated fiscal space should not be booked as the short-run funding source for large or universal social-protection commitments. Such programs would still require ordinary taxes, contributions, expenditure reallocation, or another durable source. This is not a recommendation to prioritize adoption over taxation, or taxation over adoption. Because the shortfall is joint, improving either margin alone does not close the evaluated gap. The relevant policy task is to strengthen both while keeping the program's financing plan independent of gains that have not yet materialized.

The second implication concerns sequencing by cost. Where the model finds any crossing, it appears first for a comparatively inexpensive, targeted instrument: GMI in Chile and the social pension in Peru. Starting evaluation with the least costly instrument allows a country to reach a given financing threshold earlier than beginning with a broad universal commitment. That principle is not a welfare ranking. Benefit adequacy, coverage gaps, take-up, and administrative capacity remain part of program choice. It is instead a practical ordering for fiscal appraisal: price each instrument on a common basis, examine the lower-cost options first, and expand coverage only with a financing source that remains credible beyond the favorable scenario.

The debt guardrail should then operate as a screening criterion, not as an annotation added after an accounting crossing. Peru shows why. A proposal should advance from scenario coverage to fiscal consideration only if the same revenue flow can satisfy the required debt-consistent balance. Applying that screen early prevents a temporary or gross revenue gain from being presented as free space for a permanent program. It also prevents debt stabilization and social protection from being financed simultaneously by the same fiscal dividend.

Persistence supplies the next screen. The derived horizons in Section~\ref{subsec:results_mc_hstar} show that feasibility depends on the shock lasting well beyond an immediate budget window, with the debt-consistent requirement especially demanding outside Chile. AI-linked financing is therefore a conditional long-horizon proposition. It may justify preparing scalable program designs or contingent expansion rules, but it does not justify placing recurrent obligations in the next cycle's baseline merely because a point scenario crosses eventually.

For ongoing decisions, governments can monitor the framework's observable components: realized domestic translation of the technological shock, effective fiscal capture net of behavioral erosion, and the debt-consistent primary-balance requirement. A threshold crossing should trigger re-evaluation, not automatic authorization. Budget authorities would still need to establish that the gain is realized, capturable, and sufficiently persistent before recognizing it as recurrent revenue. Until then, AI-related receipts are better treated as upside or as supplementary financing, with universal commitments backed by conventional resources.

The policy payoff is a decision rule: identify which instruments approach fiscal coherence, which accounting crossings fail the debt screen, and what evidence on persistence warrants reassessment. Section~\ref{sec:conclusion} closes with the scope of that rule.
